% from ega4-2-fr.tex
\oldpage[IV-2]{222}

que l'on a $U_{\boldsymbol R}(X')=X'$, puis, en désignant par $a$ le point
fermé de $X$, que $g'^{-1}(a)\subset U_{\boldsymbol R}(X')$. Or, soit
$x'\in g'^{-1}(a)$, et posons $z'=f'(x')$, de sorte que $b=g(z')$ est le
point fermé de $Y$. On a donc
\[
f'^{-1}(z')=f^{-1}(b)\otimes_k k(z'),
\]
et comme $k(z')$ est une extension de type fini de $k$ (puisque $g$ est
propre, donc de type fini), l'hypothèse que $\boldsymbol P(f^{-1}(b),k)$ est
vraie entraîne en vertu de la condition $4^\circ$ de (7.9.7), qu'il en est de
même de $\boldsymbol P(f'^{-1}(z'),k(z'))$. Mais comme
$\mathcal O_{Y',z'}$ est un anneau régulier et que
$\operatorname{Spec}(\mathcal O_{X',x'})\in\mathcal C$ en vertu de la
condition $1^\circ$ de (7.9.7), le lemme (7.5.1.1) prouve que
$\boldsymbol R(\mathcal O_{X',x'})$ est vraie, ce qui achève la démonstration.
\end{proof}

\begin{corollary}[7.9.9]
\label{IV.2.s7.9.9-fr}
Soient $A$, $B$ deux anneaux locaux noethériens, $\varphi:A\to B$ un
homomorphisme local, faisant de $B$ un $A$-module plat. Supposons que $A$ soit
intègre et géométriquement unibranche, et que pour toute $A$-algèbre finie
intègre $A'$ on puisse résoudre $\operatorname{Spec}(A')$ (ce qui sera par
exemple le cas si les fibres formelles de $A$ sont géométriquement régulières
et le corps résiduel de $A$ de caractéristique $0$, en raison des résultats de
Hironaka [35]). Soit $k$ le corps résiduel de $A$ et supposons que
$\operatorname{Spec}(B\otimes_Ak)$ soit géométriquement ponctuellement intègre
(4.6.9). Alors $B$ est intègre.
\end{corollary}

\begin{proof}
Nous allons appliquer (7.9.8) en prenant pour $\mathcal C$ la catégorie de
tous les préschémas localement noethériens, pour $\boldsymbol R$ la propriété
d'être intègre; les raisonnements de (7.5.3) et (7.9.7) montrent que les
conditions $1^\circ$, $2^\circ$, $3^\circ$ et $4^\circ$ de (7.9.7) sont
satisfaites dans ce cas. L'hypothèse sur $B\otimes_Ak$ et la proposition
(7.9.8) montrent donc (avec les notations de (7.9.8)) que les fibres
$f^{-1}(y)$ sont géométriquement ponctuellement intègres pour $y\in Y$;
\emph{a fortiori} les $f^{-1}(y)$ pour $y\in Y$ sont des préschémas réduits,
et comme $A$ est intègre (et \emph{a fortiori} réduit), on voit déjà que $B$
est réduit (3.3.5). Il reste à prouver que $X=\operatorname{Spec}(B)$ est
irréductible. Or, la démonstration de (7.9.8) prouve que $X'$ est localement
intègre (étant localement noethérien). Comme le morphisme $g':X'\to X$ est
surjectif, il suffit donc de voir que $X'$ est irréductible, ou (puisque $X'$
est localement intègre), que $X'$ est connexe. Mais il suffit pour cela de
prouver que la fibre $g'^{-1}(a)$ est connexe. En effet, si ce point est
établi, $X'$ ne peut être somme de deux préschémas ouverts non vides $X'_1$,
$X'_2$, car on aurait alors par exemple
$g'^{-1}(a)\cap X'_2=\varnothing$; mais la restriction $X'_2\to X$ de $g'$
étant propre (puisque $X'_2$ est un sous-préschéma fermé de $X'$),
$g'(X'_2)$ serait une partie fermée non vide de $X$ ne contenant pas le point
fermé $a$, ce qui est absurde.

Or, on a
$g'^{-1}(a)=g^{-1}(b)\otimes_{k(b)}k(a)$; mais $g:Y'\to Y$ est propre et
birationnel, donc, dans la factorisation de Stein
$Y'\xrightarrow{v}Y''\xrightarrow{u}Y$ du morphisme $g$ (\textbf{III}, 4.3.3),
le morphisme fini $u$ est aussi birationnel et par suite $Y''$ est intègre.
Comme $A$ est supposé géométriquement unibranche, il résulte alors de
(\textbf{III}, 4.3.4) que $g^{-1}(b)$ est géométriquement connexe, ce qui
achève la démonstration.
\end{proof}

\begin{remarks}[7.9.10]
\label{IV.2.s7.9.10-fr}
(i) Nous montrerons plus loin (18.8.11) que pour qu'un anneau local $B$ réduit
soit géométriquement unibranche, il faut et il suffit que pour tout morphisme
étale $h:X'\to X=\operatorname{Spec}(B)$ et tout point $x'\in X'$ au-dessus
du point fermé de $X$, $\mathcal O_{X',x'}$ soit intègre. On en déduit que si,
dans (7.9.9) on suppose, non seulement que
$\operatorname{Spec}(B\otimes_Ak)$ soit géométriquement ponctuellement
intègre, mais géométriquement unibranche, alors on peut en conclure que $B$ est
géométriquement unibranche.

